% why-anything.tex — Section 2: Why Anything?

\section{Why Anything?}
\label{sec:why-anything}

\subsection{``Physical'' is an empty word}
\label{subsec:physical-empty}

The conventional view adds a distinction to the ontology: some
mathematical structures are ``physical,'' and the rest are not.  A
particular $\SU{3} \times \SU{2} \times \UU{1}$ gauge theory with
specific coupling constants is declared to exist in a way that other
consistent structures do not.  What does ``physical'' add?

Consider the Mandelbrot set.  Its boundary is infinitely detailed, and
its patterns exist whether or not anyone computes them.  The statement
``the Mandelbrot set has a cardioid-shaped main body'' is true
independently of any substrate --- no computer needs to run, no
mathematician needs to look.  Now consider a Turing machine whose
execution trace, suitably interpreted, contains patterns equivalent to
a neural network processing sensory data.  The patterns in that trace
exist whether or not anyone runs the machine, for the same reason the
Mandelbrot set's patterns exist.

To claim that one of these patterns is ``physical'' and the other is
not, you must specify what ``physical'' adds beyond ``consistent
mathematical structure containing patterns that, from the inside,
function as observers.''  Primitivist responses exist: Maudlin argues
that ``physical'' is a primitive not requiring further
specification; Lewis's natural properties framework offers a
non-trivial account of the physical/mathematical
boundary.  We acknowledge these positions but maintain that radical
relativity posits \emph{less}: the primitivist adds an unanalyzed
distinction, while we refuse to add one.  The burden falls on the
primitivist to explain what work the distinction does beyond labeling
the structure we happen to inhabit.

\begin{postulate}[Radical Relativity]
\label{post:radical-relativity}
No instantiation of a consistent mathematical structure is
ontologically privileged over any other.
\end{postulate}

This is not an axiom being added.  It is the refusal to add the axiom
``some instantiations are special.''  The simpler position --- the one
with fewer ontological commitments --- is radical relativity.  The
burden of proof falls on anyone who would draw the line.

\begin{remark}
The parallel to Einstein is suggestive.  In 1905, Einstein did not add
a principle to electrodynamics.  He \emph{refused} to add the
assumption that one inertial frame is privileged.  The consequences
(length contraction, time dilation, $E = mc^2$) followed from the
removal.  Radical relativity makes the same move one level up:
refuse to privilege any instantiation of a mathematical structure,
and the consequences (quantum mechanics, the Standard Model) follow
from the removal.
\end{remark}

\textbf{Contrast with Tegmark.}  Tegmark's Mathematical Universe
Hypothesis~\citep{tegmark2008} asserts that ``the physical world IS a
mathematical structure.''  This \emph{elevates} mathematics to the
status of the physical --- it adds the claim that one specific
structure is ``the'' physical world.  Our move is the opposite:
\emph{deflation}.  We do not say that mathematics rises to meet
physics.  We say that ``physical'' was never doing any work.  There is
no promotion.  There is a demotion of the word ``physical'' from
ontological category to empty label.

The distinction matters.  Tegmark's framework inherits a measure
problem: if one structure is ``the'' physical world, which one?  If
all are equally physical, how do you count them?  Radical relativity
sidesteps this by refusing to designate any structure as special.  The
measure problem is replaced by a different question: given that all
structures exist, what do observers inside self-modeling structures
experience?  That question has a precise answer
(Sec.~\ref{sec:why-qm}).

\subsection{Structure space}
\label{subsec:structure-space}

\begin{definition}[Structure Space]
\label{def:structure-space}
Let $\mathcal{S}$ be the class of all consistent mathematical
structures, quotiented by isomorphism.  Two structures related by an
isomorphism are the same point in $\mathcal{S}$.  Not equivalent.
Identical.
\end{definition}

$\mathcal{S}$ is naturally a groupoid: morphisms are isomorphisms
between structures, and every morphism is invertible.  Its key
properties follow immediately from the quotient:

\begin{enumerate}[nosep]
\item \textbf{Copies collapse.}  A billion instantiations of the same
  pattern in different substrates constitute one point in $\mathcal{S}$.
  There is no counting problem, because there is nothing to count.

\item \textbf{The measure problem dissolves.}  Every multiverse
  framework struggles with the question ``how do you weight infinite
  copies?''  In $\mathcal{S}$, isomorphic copies are one thing.  The
  problematic infinity never arises.

\item \textbf{Substrate independence is automatic.}  The pattern IS the
  entity.  ``Which instantiation are you?'' is not a meaningful
  question, just as ``which coordinate system are you really in?'' is
  not meaningful after Einstein.
\end{enumerate}

\begin{remark}
Mueller's algorithmic idealism~\citep{mueller2024} provides a
natural candidate for a measure on $\mathcal{S}$: a Solomonoff-type
prior that weights simpler structures more heavily.  Our framework
is compatible with Mueller's measure but does not require it.  The
derivation chain (Table~\ref{tab:chain}) works for \emph{any} measure
that does not assign zero weight to self-modeling structures.
\end{remark}

\subsection{The relativity principle}
\label{subsec:relativity-principle}

\begin{postulate}[Level IV, restated]
\label{post:l4}
All consistent mathematical structures exist.
\end{postulate}

Under radical relativity (Postulate~\ref{post:radical-relativity}),
this is not an additional claim.  It is the same claim restated.  If
no structure is privileged, then declaring any structure nonexistent
would require specifying what distinguishes the existent from the
nonexistent --- i.e., drawing the line that radical relativity refuses
to draw.

The logical structure is:
\begin{enumerate}[nosep]
\item Radical relativity: no instantiation is privileged.
\item Therefore no structure can be declared nonexistent without
  privileging the existent ones.
\item Therefore all consistent structures exist.
\end{enumerate}

This is the sense in which Level~IV is a tautology under radical
relativity rather than an independent axiom.  It posits nothing.  It
is what remains when you stop positing.

\subsection{Idealism-structuralism convergence}
\label{subsec:idealism}

The derivation chain can be read in two directions.  From the
structuralist direction: all structures exist (Postulate~\ref{post:l4}),
the experiential measure $\rho$ selects self-modelers, and experience
lives where $\rho > 0$.  From the idealist direction: experience is
fundamental; asking what structure it has leads to self-modeling Jordan
algebras, $\hthree$, and the Standard Model gauge group.

These are not competing frameworks.  They are the same picture viewed
from two sides.  Science operates inside the region where $\rho > 0$:
every experiment, observation, and measurement occurs within a
self-modeling structure.  From inside, consciousness appears
fundamental (the idealist's observation).  From outside --- the
Level~IV view --- all structures exist and $\rho$ is one measure among
many (the structuralist's observation).  Both observations are
correct.  Neither is complete without the other.
